\section{Validation}

\subsection{RQ1: Comparative Toxicity Analysis of LLMs}

\begin{itemize}
    \item \textbf{MUST: Quantitative Validation}
    This must is completed when all the outputs, given the prompts of the dataset [...], of the three LLMs are compared and the consistency of their outputs is scored.

    \item\textbf{SHOULD: }
    This should is completed when there is a list of words that lead to a higher toxicity score. And when we can explain what kind of words these are. 
    
    \item \textbf{COULD: Qualitative Validation}
    This could is completed when the findings on why the outputs are given a high toxicity score are represented in a guideline. This guideline enables humans to (partly) prevent toxic outputs or helps lower the toxicity of outputs by showing what causes toxicity in the input.
\end{itemize}

\subsection{RQ2: Lexical Analysis of Toxic Inputs}

\begin{itemize}
    \item \textbf{MUST: Toxic Output Subset}
    If a threshold for the toxicity scores of the outputs is established this must is completed. Since there will be a trade-off between the size of the usable dataset and the height of the toxicity score, a graph will be used to identify the preferred threshold.
    
    \item \textbf{MUST: Attribution Explanation}
    This must is completed if it is possible to visually show which inputs contribute most to generating toxic outputs.
    
    \item \textbf{MUST: Lexical Analysis}
    This is completed if there is a lexical understanding of the tokens that lead to toxicity. Is it mainly foul language, is it a sensitive subject, etc. 

    \item \textbf{COULD: Human Evaluation Protocol}
    This is completed if there are relevant findings that are added to the previously defined guideline. It may be possible that there are no relevant findings that can be added.
\end{itemize}

\subsection{RQ3: Syntactic Analysis of Toxic Inputs}

\begin{itemize}
    \item \textbf{SHOULD: Syntactic Analysis}
    This should is completed when parsing methods are used to analyse the structure of toxic inputs and an accuracy is scored.
    
    \item \textbf{SHOULD: Integration with Lexical Analysis}
    This should is completed when the syntactic and lexical findings are combined and represented in a guideline.
\end{itemize}
